% !TEX root = ../thesis.tex
\chapter*{MỞ ĐẦU}
\addcontentsline{toc}{chapter}{MỞ ĐẦU}

\section*{Tính cấp thiết của đề tài}

Tin tức trực tuyến tiếng Việt đã trở thành phương tiện truyền thông thông tin công cộng
chủ yếu, với hàng chục triệu bài báo được đăng tải mỗi năm trên các tờ báo lớn như
\textit{tuoitre.vn}, \textit{thanhnien.vn}, \textit{tienphong.vn},
\textit{vietnamnet.vn} và nhiều tờ khác. Khối lượng thông tin khổng lồ này tạo ra hai
nhu cầu song song cho người đọc. Nhu cầu đầu tiên là \textbf{nén thông tin}: một bản tóm
tắt súc tích, đáng tin cậy giúp người đọc có thể nắm bắt nội dung bài báo dài chỉ trong
vài giây. Nhu cầu thứ hai là \textbf{xác thực thông tin}: kiểm tra xem các tuyên bố chính
trong bài báo có được các nguồn uy tín khác xác nhận hay không, đặc biệt trong bối cảnh
tin giả, tiêu đề giật gân và các bài viết chỉ phản ánh một phần sự thật lan truyền song
song cùng các bài báo chính xác.

Hiện nay, các hệ thống tóm tắt tiếng Việt xây dựng dựa trên một Mô hình ngôn ngữ lớn
(LLM) đơn lẻ vẫn cho ra kết quả không ổn định --- đôi khi quá dài, đôi khi quá sát từ
ngữ gốc, thậm chí đôi khi bịa đặt thông tin. Một hướng nghiên cứu gần đây,
\textbf{Mixture-of-Agents} (MoA) của Wang et al.\ \cite{Wang2024}, lập luận rằng
nhiều LLM có thể được kết hợp trong một kiến trúc phân cấp tác nhân đề xuất--tổng hợp
để tạo ra kết quả vượt trội so với bất kỳ mô hình đơn lẻ nào trên các thang đánh giá.
Tuy nhiên, liệu phương pháp này có thể áp dụng với bài toán tóm tắt tiếng Việt
\textit{có neo ngữ cảnh gốc} --- nơi hệ thống phải trung thực với bài báo nguồn thay
vì trả lời tự do --- hay không, vẫn là câu hỏi chưa được giải đáp trong các nghiên
cứu xử lý ngôn ngữ tự nhiên tiếng Việt.

Một vấn đề sâu hơn thúc đẩy khóa luận này: các hệ số đo lường phổ biến nhất trong
nghiên cứu tóm tắt tiếng Việt --- ROUGE \cite{Lin2004}, BLEU và BERTScore
\cite{Zhang2020} được tính toán dựa vào đối chiếu so với bài báo gốc --- đo độ trùng lặp, không đo chất lượng.
Các thang điểm này đánh giá cao các bản tóm tắt tái sử dụng cụm từ từ bài gốc và chấm điểm thấp các bản tóm tắt
diễn giải lại hoặc rút gọn mạnh. Một tác nhân tổng hợp được sử dụng đúng, tổng
hợp các bản thảo thành bản tóm tắt biên tập sạch hơn, về nguyên lý sẽ đạt điểm thấp
hơn trên các hệ số này so với bản thảo sao chép cụm từ đơn thuần. Nếu nghiên cứu tóm
tắt tiếng Việt muốn áp dụng các hệ thống kiểu tác nhân tổng hợp, lĩnh vực này cần
một phương pháp luận đánh giá không có thiên kiến cấu trúc bất lợi cho các hệ thống đó.

\section*{Ý nghĩa khoa học và thực tiễn}

\textbf{Ý nghĩa khoa học.}
Khóa luận đóng góp một \textit{khung đánh giá ba trục} kết hợp: (i)~các hệ số overlap
nội dung so với bài gốc, (ii)~đánh giá chất lượng bởi LLM và xếp hạng theo
cặp, và (iii)~xếp hạng mù bởi con người với chỉ số đồng thuận giữa các người đánh giá.
Khóa luận sử dụng khung đó để chỉ ra, bằng bằng chứng thực nghiệm trên dữ liệu tiếng
Việt, một hạn chế cấu trúc của các hệ số overlap đối với bài toán tóm tắt kiểu tổng
hợp. Hạn chế này đã được đề cập gián tiếp trong các nghiên cứu tiếng Anh
\cite{Wang2024,Zheng2023} nhưng chưa được nghiên cứu tường minh cho tiếng Việt.

\textbf{Ý nghĩa thực tiễn.}
Khóa luận cung cấp một sản phẩm đầu-cuối hoàn chỉnh có tên \textbf{Fiber}, mà bất
kỳ người đọc báo Việt Nam nào cũng có thể cài đặt trên trình duyệt dựa trên Chromium.
Fiber nhúng một thanh sidebar vào bảy tên miền báo chí tiếng Việt lớn, tóm tắt bài
báo theo thời gian thực bằng một LLM do người dùng chọn (hoặc pipeline MoA dung hợp),
và kiểm chứng bản tóm tắt so với bằng chứng từ web trực tiếp thông qua API tìm kiếm
Tavily. Toàn bộ hệ thống --- một backend Next.js, một tiện ích mở rộng trình duyệt
Plasmo, và một dịch vụ BERTScore Python --- là mã nguồn mở.

\section*{Mục tiêu nghiên cứu}

\begin{enumerate}
  \item Thiết kế và xây dựng hệ thống tiện ích mở rộng trình duyệt đầu-cuối hoàn
        chỉnh, có khả năng tóm tắt và kiểm chứng bài báo tiếng Việt theo thời gian
        thực, hỗ trợ nhiều nhà cung cấp LLM (OpenAI, Anthropic, Google) và nhiều cơ chế
        phân loại định tuyến.
  \item Thích ứng pipeline Mixture-of-Agents của Wang et al.\ (2024) cho bài toán
        tóm tắt tiếng Việt có neo ngữ cảnh gốc, bao gồm một kết nối phần dư bài
        viết nguồn vào câu lệnh của tác nhân tổng hợp.
  \item Đề xuất và kiểm chứng một khung đánh giá ba trục --- khả năng giữ lại nội dung
        (Trục~A), chất lượng kết quả tóm tắt đánh giá bởi LLM (Trục~B), và xếp hạng
        con người với chỉ số đồng thuận (Trục~C) --- đồng thời sử dụng khung đó
        để so sánh dung hợp MoA với các mô hình đơn lẻ cơ sở trên tập dữ liệu
        tiếng Việt cân bằng chủ đề.
\end{enumerate}

\section*{Đối tượng và phạm vi nghiên cứu}

\textbf{Đối tượng.} Các bài báo tiếng Việt được đăng trên
\textit{tuoitre.vn}, \textit{thanhnien.vn}, \textit{tienphong.vn},
\textit{vietnamnet.vn}, \textit{laodong.vn}, \textit{vtv.vn} và
\textit{nld.com.vn}. Thực nghiệm đánh giá trong Chương~4 sử dụng 154 bài báo
từ \textit{tienphong.vn}, được cân bằng theo năm chủ đề: thời sự, pháp luật,
kinh tế, giáo dục và văn hóa.

\textbf{Phạm vi.}
\begin{itemize}
  \item \textit{Mô hình.} Ba LLM thương mại trên đám mây (GPT-4o-mini, Claude
        Haiku 4.5, Gemini 2.5 Flash) làm tác nhân đề xuất; GPT-4o làm tác nhân
        tổng hợp; lựa chọn bản thảo tốt nhất làm mô hình đơn lẻ cơ sở mạnh nhất.
        Các mô hình đặc thù tiếng Việt (ViT5 \cite{Nguyen2022} và
        PhoBERT \cite{Nguyen2020}) được sử dụng cho phân loại định tuyến và
        đánh giá dựa trên nhúng từ.
  \item \textit{Đánh giá.} Đơn tên miền (\textit{tienphong.vn}), một lượt chạy
        mỗi bài báo, không có yêu cầu tổng quát hóa đa tên miền. Nghiên cứu con
        người mang tính bao quát với $n = 17$ người đánh giá và $48$ tác vụ, được
        đóng khung như một bổ sung mô tả cho Trục~A và Trục~B.
  \item \textit{Kiểm chứng thông tin.} Xác minh tăng cường tìm kiếm qua API tìm
        kiếm web Tavily. Phát hiện tuyên bố và truy xuất bằng chứng đã được cài
        đặt; kiểm chứng thông tin có cấu trúc sâu hơn (đồ thị tri thức hoặc phân
        loại quan điểm) nằm ngoài phạm vi của khóa luận này.
\end{itemize}

\section*{Phương pháp nghiên cứu}

Khóa luận áp dụng phương pháp thực nghiệm so sánh. Cùng một tập 154 bài báo được tóm
tắt bởi năm phương pháp (ba tác nhân đề xuất riêng lẻ, đầu ra dung hợp MoA, và mô hình
đơn lẻ cơ sở lựa chọn bản thảo tốt nhất), và 770 bản tóm tắt thu được được đánh giá
theo ba trục. Trục~A sử dụng ROUGE-1/2/L, BLEU, BERTScore (với PhoBERT) và tỷ lệ nén;
Trục~B sử dụng rubric 5 chiều theo FLASK \cite{Ye2023}, điểm tuyệt đối kiểu MT-Bench
\cite{Zheng2023}, đánh giá sở thích từng cặp kiểu AlpacaEval có ngẫu nhiên hóa vị trí
\cite{Dubois2024}, và bộ phân loại tính xác thực cấp tuyên bố; Trục~C sử dụng giao
diện xếp hạng mù $K = 3$ được quản lý thông qua chính sản phẩm tiện ích mở rộng trình
duyệt. Các kiểm định thống kê bao gồm kiểm định dấu hai chiều (loại trừ hòa) cho tỷ
lệ thắng từng cặp và Fleiss' $\kappa$ \cite{Landis1977} cho đồng thuận giữa các người
đánh giá; thiên kiến độ dài được kiểm soát bằng phương pháp phân nhóm đơn giản hóa của
Dubois et al.\ \cite{Dubois2024}.

\section*{Bố cục khóa luận}

Phần còn lại của khóa luận được tổ chức như sau.

\textbf{Chương 1 --- Tổng quan.} Giới thiệu bài toán tóm tắt và kiểm chứng tin tức,
kỹ thuật Mixture-of-Agents, và các công trình liên quan trong nghiên cứu tiếng Việt
cũng như quốc tế, từ đó xác định khoảng trống mà khoá luận sẽ khai thác.

\textbf{Chương 2 --- Cơ sở lý thuyết.} Trình bày các kiến trúc LLM được sử dụng
trong hệ thống, ba dòng hệ số đánh giá (overlap, LLM-giám khảo, xếp hạng con người),
phương pháp luận tính xác thực, và các công nghệ để xây dựng tiện ích mở rộng trình
duyệt cùng các dịch vụ backend.

\textbf{Chương 3 --- Thiết kế và xây dựng hệ thống.} Trình bày phân tích yêu cầu,
kiến trúc ba dịch vụ, lớp trừu tượng hóa LLM đa nhà cung cấp, pipeline MoA (đóng
góp kỹ thuật), khung đánh giá ba trục (đóng góp phương pháp luận), lược đồ cơ sở
dữ liệu, và giao diện người dùng.

\textbf{Chương 4 --- Thực nghiệm và đánh giá.} Báo cáo kết quả theo ba trục trên
tập dữ liệu 154 bài báo, bao gồm kết quả từng cặp nổi bật nhất (dung hợp thắng bản
thảo tốt nhất $68,9\%$, $p < 0,0001$), phân tích theo chủ đề, diễn giải nghiên cứu
con người, tam giác hóa đa trục và các hạn chế phương pháp luận.

\textbf{Kết luận} tóm tắt các đóng góp, hạn chế và hướng phát triển tương lai;
ba phụ lục trình bày ví dụ minh họa, thử nghiệm lời nhắc tác nhân tổng hợp và
pipeline kiểm chứng thông tin.
